\subsection{Future Directions}
\label{appendix:fw}

Of the error types we catalog in our analysis of \texttt{LINC}, the key opportunity for improvement is more elegant handling of naturalistic language use.
While the errors observed with \texttt{CoT} result from faulty deductive inferences, in the case of \texttt{LINC}, all errors have been localized to the semantic parsing procedure.
This process flows primarily unimpeded in the evaluation of the synthetic ProofWriter dataset, yet leaves room for improvement with the naturalistic FOLIO. 
In follow-up work, we hope to deeply explore naturalistic evaluation settings, as when data get the most messy is also where improvements become the most valuable. 
Here, we propose three strategies for further improvement on naturalistic settings.

First, in naturalistic communication, ``obvious'' information is often left out of explicit productions, left to be inferred in the ``common ground'' of the communicative act \citep{grice1975logic,stalnaker2002common}. 
Implicit premise rediscovery through controlled exploration on the logical neighborhood of the existing explicit premises promises to be a powerful strategy for improving performance in underspecified settings.

Second, while a number of samples are lost to syntax errors, recent work has proposed restricting the sampling space of an LLM to that which is consistent with term expansions in a context-free-grammar (CFG) \citep{poesia2022synchromesh}.
Doing so in this setting would eliminate all syntax errors.

Third, sometimes, the translation process to FOL is lossy, throwing away valuable information present in the original sentence.
We propose improving the faithfulness of FOL translations by asking the LLM to translate the FOL back to natural language and comparing with the original.
Forward translations that rank highly when back-translated would be those which have effectively captured the intricacies of a particular sentence's semantics.

Overall, we believe that shifting to evaluations on more naturalistic datasets, and incorporating strategies such as those presented here, will help pave the path forward for neurosymbolic approaches to formal reasoning.